% ============================================================
% harryopo-tikz-diagram - 组织架构图模板 (Organization Tree)
% ============================================================
%
% 【用途】
%   绘制树形组织架构图，支持 2-3 层级，每个节点包含职位和姓名。
%   使用 TikZ 的 trees 库，连接线为直角分叉。
%
% 【使用方法】
%   将本文件内容复制到你的 LaTeX 文档中，或使用 \input 命令引入。
%   修改下方"可调整参数"部分来自定义你的组织架构图。
%
% 【编译要求】
%   XeLaTeX 编译器，需加载 ctex 宏包以支持中文。
%
% ============================================================

% ------------------------------------------------------------
% 必要的宏包（如果文档中已加载可省略）
% ------------------------------------------------------------
% \usepackage{tikz}
% \usetikzlibrary{positioning, fit, backgrounds, arrows.meta}
% \usetikzlibrary{trees}

% ------------------------------------------------------------
% 可调整参数 —— 在这里修改你的组织架构图配置
% ------------------------------------------------------------

% 主题色定义（Nature 风格）
\definecolor{PaperBorder}{HTML}{B0C4DE}      % 柔和灰蓝边框
\definecolor{PaperFill}{HTML}{FFFFFF}        % 纯白节点填充
\definecolor{PaperTitle}{HTML}{2C3E50}       % 深蓝灰标题
\definecolor{PaperLayerBg}{HTML}{F5F8FA}     % 极淡蓝灰背景
\definecolor{PaperLayerBorder}{HTML}{D8E3EE} % 层边框：更浅的灰蓝
\definecolor{PaperMuted}{HTML}{5D6D7E}       % 次要文字：灰蓝
\definecolor{AccentOrange}{HTML}{D35400}     % 强调色：柔和橙
\definecolor{ShadowColor}{HTML}{E0E8F0}      % 阴影色：极淡蓝灰

% 树的方向
%   down  = 自上而下（默认，根在顶部）
%   up    = 自下而上（根在底部）
%   right = 从左到右（根在左侧）
%   left  = 从右到左（根在右侧）
\def\orgTreeGrow{down}

% 层级间距
\def\orgLevelOneSep{60pt}     % 第一层到第二层的距离
\def\orgLevelTwoSep{50pt}     % 第二层到第三层的距离
\def\orgSiblingOneSep{20pt}   % 第一层子节点间距
\def\orgSiblingTwoSep{15pt}   % 第二层子节点间距

% 节点尺寸
\def\orgNodeWidth{90pt}       % 节点宽度（text width）
\def\orgNodeHeight{50pt}      % 节点最小高度
\def\orgRootWidth{120pt}      % 根节点宽度

% 连接线样式
\def\orgLineThickness{1.2pt}  % 连接线粗细
\def\orgLineColor{PaperBorder}  % 连接线颜色

% ------------------------------------------------------------
% 样式定义（一般不需要修改）
% ------------------------------------------------------------
\tikzset{
  % ===== 根节点样式（Nature 风格） =====
  org-root/.style = {
    draw = PaperBorder,
    fill = PaperTitle,
    text = PaperFill,
    line width = 0.8pt,
    rounded corners = 6pt,
    text width = \orgRootWidth,
    minimum height = \orgNodeHeight,
    align = center,
    font = \small\bfseries\sffamily,
    inner sep = 6pt,
  },
  % ===== 二级节点样式 =====
  org-level2/.style = {
    draw = PaperBorder,
    fill = PaperFill,
    line width = 0.8pt,
    rounded corners = 4pt,
    text width = \orgNodeWidth,
    minimum height = \orgNodeHeight,
    align = center,
    font = \small\sffamily,
    text = PaperTitle,
    inner sep = 5pt,
  },
  % ===== 三级节点样式 =====
  org-level3/.style = {
    draw = PaperLayerBorder,
    fill = PaperFill,
    line width = 0.6pt,
    rounded corners = 3pt,
    text width = \orgNodeWidth,
    minimum height = \orgNodeHeight,
    align = center,
    font = \footnotesize\sffamily,
    text = PaperTitle,
    inner sep = 4pt,
  },
  % ===== 部门组背景框 =====
  org-groupbox/.style = {
    draw = PaperLayerBorder,
    fill = PaperLayerBg,
    dashed,
    line width = 0.6pt,
    rounded corners = 6pt,
    inner sep = 10pt,
  },
  % ===== 连接线 =====
  org-edge/.style = {
    draw = PaperBorder,
    line width = \orgLineThickness,
  },
}

% ------------------------------------------------------------
% 模板正文 —— tikzpicture 环境
% ------------------------------------------------------------
%
% 【使用说明】
%   1. 使用 TikZ 的 tree 语法：child { node { ... } child { ... } }
%   2. 根节点放在 (0,0) 位置
%   3. 节点内容用 \textbf{职位} \\ 姓名 格式
%   4. 使用 edge from parent fork down 实现直角分叉
%
% 【节点内容格式】
%   每个节点包含两行：
%   第一行：职位（加粗）
%   第二行：姓名（常规）
%
% ------------------------------------------------------------

\begin{tikzpicture}[
    % 树的全局设置
    grow = \orgTreeGrow,
    edge from parent fork down,
    % 第一层子节点
    level 1/.style = {
      sibling distance = \orgSiblingOneSep,
      level distance = \orgLevelOneSep,
      edge from parent/.style = {org-edge},
      nodes = {org-level2},
    },
    % 第二层子节点
    level 2/.style = {
      sibling distance = \orgSiblingTwoSep,
      level distance = \orgLevelTwoSep,
      edge from parent/.style = {org-edge},
      nodes = {org-level3},
    },
    % 所有节点的默认设置
    every node/.style = {inner sep = 0pt, outer sep = 0pt},
  ]

  % ==========================================================
  % 示例：软件公司团队组织结构
  % ==========================================================

  % 根节点：总经理（第一个节点用绝对坐标 at (0,0)）
  \node[org-root] at (0,0) {
    \textbf{总经理} \\
    张总
  }
  % ===== 一级部门：技术部 =====
  child {
    node {
      \textbf{技术总监} \\
      李工
    }
    child { node {\textbf{前端组} \\ 王小明} }
    child { node {\textbf{后端组} \\ 陈大华} }
    child { node {\textbf{测试组} \\ 刘芳} }
    child { node {\textbf{运维组} \\ 赵强} }
  }
  % ===== 一级部门：产品部 =====
  child {
    node {
      \textbf{产品总监} \\
      孙敏
    }
    child { node {\textbf{产品设计} \\ 周杰} }
    child { node {\textbf{UI 设计} \\ 吴丽} }
  }
  % ===== 一级部门：运营部 =====
  child {
    node {
      \textbf{运营总监} \\
      郑华
    }
    child { node {\textbf{市场运营} \\ 钱伟} }
    child { node {\textbf{用户运营} \\ 冯静} }
  }
  % ===== 一级部门：行政部 =====
  child {
    node {
      \textbf{行政总监} \\
      马琳
    }
    child { node {\textbf{人事行政} \\ 朱琳} }
    child { node {\textbf{财务} \\ 徐峰} }
  };

  % ==========================================================
  % 可选：部门分组背景框（使用 fit 库）
  % ==========================================================
  % 如需给某个部门加背景框，取消下方注释并修改节点名
  %
  % \begin{scope}[on background layer]
  %   \node[org-groupbox, fit=(tech-node1)(tech-node2)(tech-node3)] (tech-group) {};
  %   \node[below right=6pt of tech-group.north west, font=\small\bfseries, text=SubColor]
  %     {技术中心};
  % \end{scope}

\end{tikzpicture}

% ============================================================
% 【常见修改示例】
% ============================================================
%
% 1. 增加/减少层级：
%    在 child 中再嵌套 child 即可增加层级。
%    最多建议 3 层，再多会导致排版拥挤。
%
% 2. 调整子节点数量：
%    增加或删除 child { node { ... } } 块。
%    同步调整 sibling distance 以适应节点数量。
%
% 3. 修改树的生长方向：
%    将 \def\orgTreeGrow 改为 down / up / left / right
%
% 4. 自定义节点内容：
%    节点内可以是任意内容，不一定非要是 职位+姓名。
%    例如：node {部门名称} 或 node {团队代号}
%
% 5. 调整节点大小：
%    修改 \orgNodeWidth, \orgNodeHeight, \orgRootWidth
%    或在节点上单独指定：
%    node[org-level2, text width=100pt] { ... }
%
% 6. 添加部门标题：
%    使用 fit 库在背景层添加分组框和标题，见模板中的示例。
%
% 7. 改变连接方式：
%    - 直角分叉：edge from parent fork down（默认）
%    - 直线连接：移除 fork 样式，使用默认直线
%    - 曲线连接：edge from parent/.style={bend left}
%
% ============================================================
